\chapter{The economic case for a shared product}\label{ch:economics}
\section{Compare the same scientific work}
The credible alternative is a capable internal team using its warehouse,
catalog, versioned notebooks and established open-source packages. That
team can implement the methods surveyed here. The proposed product must
justify the extra expense of reusable semantics, interfaces and operations
through repeated use, lower review effort or additional decision value.
Its comparison cannot assume that the alternative omits temporal correctness,
credible controls or model validation.

Define a completed study as a decision with a frozen population and
information clock, a stated target, defensible identifying restrictions,
at least one suitable reference comparison, an economic account and a
reconstructible evidence package. Both alternatives must deliver the same
number and mix of such studies, with equivalent scientific review.
The platform automates recurring work; expert composition performs more
of it in study-specific code and procedures. Neither is entitled to count
unreviewed model fits as completed studies.

The cost account has initial common construction $K_j$, annual shared
operation $F_{jt}$ and incremental study cost $c_{jk}$ for alternative
$j$. Study cost includes analyst, engineering and review effort, allocated
compute and recurring data reconciliation. Shared cost includes dependency
qualification, orchestration, support, security and reusable validation.
An existing bank catalog is a common baseline input; a new connection
required by one alternative is an incremental cost of that alternative.

\section{The reuse threshold}
For $n_t$ equivalent completed studies in year $t$, present cost is
\begin{equation}
 C_j=K_j+\sum_{t=1}^{H}
 \frac{F_{jt}+\sum_{k=1}^{n_t}c_{jk}}{(1+r)^t}.
 \label{eq:equivalent-cost}
\end{equation}
Suppose for illustration that annual shared cost and per-study cost are
constant, with equal study volume $n$. Let
$A_H(r)=\sum_{t=1}^{H}(1+r)^{-t}$,
$\Delta K=K_P-K_E$, $\Delta F=F_P-F_E$, and
$\Delta c=c_E-c_P>0$, where $P$ denotes the product and $E$ expert
composition. Then the product has lower construction/operation cost when
\begin{equation}
 n>\frac{\Delta K/A_H(r)+\Delta F}{\Delta c}.
 \label{eq:reuse-threshold}
\end{equation}
The numerator is its annualized common-cost premium; the denominator is
the saving per completed equivalent study. If $\Delta c\leq0$ and
common-cost premiums are positive, reuse alone cannot justify the platform.
No amount of repeating an equally expensive workflow removes that premium.

Heterogeneous studies need the original sum in
Equation~\eqref{eq:equivalent-cost}. A new country's undocumented action
logs may require nearly as much work with either stack. A twentieth
deposit rollout over a qualified data account may reuse most of its
transformation and review evidence. Counting both as identical use cases
would inflate the expected saving.

\section{Illustrative costs, including losing scenarios}
Use a three-year horizon and 10\% discount rate. All amounts in this section
are illustrative USD millions. Expert composition requires initial common
work of 1.0, annual shared cost 0.6 and 0.15 per completed study.
The product requires 4.0 initially, 1.8 annually and 0.06 per study.
These are planning assumptions to be replaced by a bank estimate. They are
not observed salary rates, procurement prices or measured productivity.

The study-cost difference is 0.09 million. At the separate staffing
assumption of USD\,25,000 per person-month, that is 3.6 person-months
saved per completed study. That is a large claim: it must be demonstrated
on equivalent tasks, not inferred from the existence of a common API.
The product's annual shared cost includes support and dependency work;
study-specific engineering and review are included in its per-study cost.
The categories must not overlap.

\begin{table}[htbp]\centering\small
\begin{tabular}{@{}lrrr@{}}
\toprule Scenario & Expert cost & Platform cost & Platform saving\\\midrule
Low reuse & 5.476 & 9.670 & -4.194\\
Moderate reuse & 11.445 & 12.057 & -0.613\\
High reuse & 15.921 & 13.848 & 2.073\\
Adoption ramp & 8.402 & 10.840 & -2.438\\
\bottomrule\end{tabular}

\caption{Three-year present cost in illustrative USD millions. Constant
annual volumes are 8, 24 and 36 studies; the ramp is 4, 16 and 30.
Positive saving favors the shared product.}
\label{tab:reuse-economics}
\end{table}

The cost-only threshold is about 26.74 completed studies per year. At
eight annually the product costs about USD\,4.19 million more in present
value. At 24 it still costs about USD\,0.61 million more. At 36 it saves
about USD\,2.07 million. An adoption ramp of 4, 16 and 30 studies leaves a
USD\,2.44 million premium even though final-year throughput exceeds the
constant-volume threshold. Early costs occur before much reuse.

\begin{figure}[htbp]\centering
\includegraphics[width=.9\textwidth]{derivations/reuse_economics.pdf}
\caption{Equal-output cost comparison under declared assumptions. The
crossover is a hypothesis about repeatable work and costs, not a forecast
of bank adoption.}
\label{fig:reuse-economics}
\end{figure}

If the per-study saving is half as large, the threshold doubles to about
53.47 studies annually. If the annual shared-cost premium rises from
1.2 to 2.0 million, the original per-study saving needs about 35.63
studies annually. These sensitivities argue for measuring reusable effort
early and keeping common construction narrow until multiple teams use it.
A platform built for imagined reuse can be economically worse than a
well-maintained expert toolkit.

\section{Separate decision value from construction savings}
Let $D_t$ be additional realized business value attributable to the product
relative to the equivalent expert workflow. The incremental investment case is
\begin{equation}
 \Delta\mathrm{NPV}=C_E-C_P+
       \sum_{t=1}^{H}\frac{D_t}{(1+r)^t}.
 \label{eq:incremental-product-npv}
\end{equation}
The first term already counts workflow cost savings. Do not add the same
released analyst time again as a separate business benefit. If that time
supports additional studies, model their costs and incremental decisions
explicitly, with a capacity constraint.

At 24 studies per year the illustrative cost premium requires about
USD\,0.246 million of additional constant annual decision benefit.
For the adoption ramp it requires about USD\,0.980 million. Those may be
modest amounts for a large banking group, but group size is not evidence
that the product creates them. The expert comparator may already find the
same beneficial intervention. An analyst's higher point estimate is not
realized incremental value.

The deposit study gives a useful distinction. Discovering that a promotion
loses contribution can avoid an expansion loss, but the prevented loss is
attributable only if the comparator would otherwise expand it and the
finding changes the decision. The collections simulation's positive policy
gain is a property of its generating model; it supplies no bank revenue
for Equation~\eqref{eq:incremental-product-npv}. A hidden-state module may
have positive research value while its operational policy benefit remains
unresolved.

Risk reduction is also a counterfactual claim. If the platform catches a
late-data feature, record its frequency, review disposition and consequence
under the baseline process. Do not multiply an arbitrary low-probability
catastrophe by a large bank balance sheet and call the result a measured
benefit. Early validation can instead report concrete errors detected and
the cost of resolving them.

\section{Adoption, maturity and cash timing}
Not every initiated study reaches a usable conclusion. A simple planning
decomposition is
\[
 n_t^{\rm usable}=n_t^{\rm started}\,
 p_t^{\rm data}\,p_t^{\rm design}\,p_t^{\rm completed}
 \,p_t^{\rm review},
\]
where the probabilities are conditional progression rates, not assumed
independent events. The multiplicative form is the chain rule when each
rate conditions on prior progression. Costs incurred by failed studies
remain in the account. A well-designed product may appropriately terminate
more unidentified studies early; usable throughput should then be assessed
alongside avoided futile effort.

Decision outcomes mature later still. A six-month deposit outcome can
become observable within a pilot year; losses from a credit-limit change
may take longer. A continuing value model introduces assumptions rather
than making those outcomes observed. Discounting should follow the actual
date of expenditure and realized benefit, with distinct data-access,
development, study and outcome lags.

The earlier aggregate illustration can be retained as a stress calculation:
a USD\,4 million outlay and USD\,1.8 million annual operation cost need
constant annual incremental benefit of approximately
USD\,\NumBreakEven{} million over three years at 10\%.
At USD\,5 million annual benefit its NPV is USD\,\NumBaseNPV{} million;
losing the first year's benefit makes it about negative USD\,0.588 million.
Those values use a coarser account than the equal-output comparison and
must not be added to it. The explicit comparator model is the preferred
investment account.

\section{Measuring the product hypothesis in the first tranche}
Choose a small set of matched tasks across the three workspaces: reconstruct
a rollout, evaluate a repeated-decision policy under a defensible design,
and organize a temporal discovery dataset with known semantic ambiguities.
Agree their required outputs and reviewer rubric before timing them.
Give the expert-composition arm the same source access, compute budget and
scientific expertise. It may reuse its own sensible templates.

Measure active analyst and reviewer hours separately from calendar delay,
data-owner response time and machine time. Record task complexity,
previous familiarity, failed starts, rework and unresolved material findings.
A counterbalanced assignment of comparable tasks can reduce learning-order
effects; repeated use by several teams is necessary to estimate reuse
rather than first-author familiarity. A proposed target of 30\% lower
median analyst/reviewer effort is an investment hypothesis, not a result
already established by this proposal.

The platform should also export a study to a second analyst who did not
author it. Measure time to reproduce the target, source account and
financial conclusion, and whether the analyst notices deliberately
introduced temporal, support or sign errors. Equivalent review quality is
part of the output definition. Faster execution with more unrecognized
scientific defects is not a saving.

The first tranche is six people for six months: one product/economics lead,
two causal/econometric scientists, two data/platform engineers and one
application engineer. At the illustrative loaded rate, staffing is
USD\,900,000; compute/tooling USD\,50,000 and 20\% contingency bring the
total to USD\,1.14 million. That tranche is included within the illustrative
USD\,4 million platform construction, not added to it. Bank-owner,
independent-validation and security participation remain unpriced and
must be estimated consistently for both alternatives before funding.

Six people cannot production-qualify every method in the census in six
months. They can construct the common account and evidence layer, qualify
a narrow panel/longitudinal path, and deliver finite-state estimation and
temporal discovery as explicit research capabilities. The next investment
depends on observed scientific usefulness and reusable effort rather than
on completing a broad feature checklist.

\section{Independent product economics and procurement}
An independent offering adds customer-specific deployment, support,
upgrades, contracts and onboarding costs. Code reuse does not eliminate
each institution's semantic mapping and validation. A plausible initial
commercial form is a design-partner deployment with a paid construction
scope and recurring support, preserving SDK and evidence export.
Willingness to pay, sales cycle, support load and customer concentration
have not been measured here.

Compare a commercial vendor on the same reconstructed studies, permission
boundaries, supported identification routes, exported uncertainty and
dependency responsibilities. Vendor demos and list prices are insufficient;
no current private enterprise capability or price is asserted in this survey.
For internal adoption, chargeback should distinguish common capability
from a business unit's bespoke model so that demand is not inflated by
unpriced customization.

The sponsor's decision is therefore conditional but concrete: fund a narrow
shared construction if named teams have repeatable decisions, reconstructible
data and a credible path to reusable work. Continue expert composition where
volume is low or assumptions differ too much. Expand the common product
when measured study cost, quality and adoption support the incremental
account. The financial case becomes stronger through evidence of actual
reuse, not through assigning all of a bank's causal-analysis benefit to a
new software layer.
